\section{Experiments}
\label{sec:experiments}

\textit{Metric focus.}
For model training we decided to train our model over several epochs since the facial expressions cannot be learned reliably from a single pass over the data. After each batch the model parameters will be updated using backpropagation. Alongside the training rate, a validation split was used to get a better idea of how well the model generalises to unseen samples. This is especially important facial emotion recognition, where models tend to learn dataset specific patterns instead of expression related features which leads to overfitting. After each epoch, the training accuracy was calculated using our accuracy function alongside the training loss by comparing the predicted labels with the true labels.  

The Adam optimiser was chosen as our main optimisation method, as it can adjust parameter updates automatically and is generally the most common choice in deep learning. The initial tests were conducted to figure out the most optimum learning rate for Adam, and the model was trained several different times over many epochs with different learning rates between \( 0.0003 \) and \( 0.003 \).

To limit overfitting, several regulations methods were applied. Our main method to limit overfitting was Dropout to prevent the model from relying too strongly on individual neurons. In addition, after witnessing overfitting in the our initial tests \( L_2 \) weight decay was applied to keep the model parameters from growing too large. 

As well as the two regularization methods our model was tested with different epoch counts to figure out the sweet spot between underfitting and overfitting. Over \( 15 \) combinations between these three parameters were tested several times to find the average values for each combination, in order to figure out the parameter values that give the best test accuracy. The dropout rates between \( 0.1 \) and \( 0.5 \) were evaluated. To determine the best weight decay rate for \( L_2 \) regularization, tests were carried out between three different weight decay rates of \( 10^{-3} \), \( 10^{-4} \) and \( 10^{-5} \). 

Different epoch counts were also tested as we witnessed that the initial approach of training our model for \( 10 \) epochs didn’t have the perfect training accuracy scores resulting in underfitting. Our parameters were tested with \( 20 \) and then \( 30 \) epochs as more than \( 30 \) epochs produced training accuracy rates of more than \( 99\% \) which indicates overfitting. As a final parameter the batch size was assessed as we tried out the different batch sizes of \( 64 \), \( 128 \) and \( 256 \).
After implementing the custom transformation, we noticed that \( 30 \) epochs was not sufficient and that we needed more epochs to prevent underfitting. We also examined that the loss decreased slower after implementing our custom transformation but that it reached the same level after more epochs which was another reason behind training our model for more epochs. At this stage of testing we also implemented our evaluation functions to calculate the precision, recall and \( F_1 \)-score of our model.

Furthermore we implemented our landmark model, alongside our new dataset. Noticing the overfitting we decided to implement further functions to improve the test accuracy of our model. These functions include a Sobel feature branch which creates a separate CNN branch for edge information, and a width multiplier support to make experimentation with different image resolutions easier, which affects regularization implicitly.
